% Figure 1 — the germ identity, drawn to be broadly readable (not heavily formal).
% Requires: tikz (\usetikzlibrary{arrows.meta,decorations.pathreplacing,calc,positioning}).
% Placement [!ht], not [t]: main.tex inputs this file directly after the introduction's
% breakable claims box, and a [t] float takes the top of the page that box is still running
% down -- landing between its two halves. [!ht] sets the figure after the box instead.
\begin{figure}[!ht]
\centering
\begin{tikzpicture}[x=1cm,y=1cm,font=\small,>={Stealth[length=2mm]}]
  % input
  \node[draw,rounded corners,minimum width=1.6cm,minimum height=1.2cm,fill=gray!8] (x) at (0,0) {$x$};
  \node[below=1mm of x,font=\scriptsize,text width=2.2cm,align=center] {input\\($d$ coordinates)};
  % network
  \node[draw,rounded corners,minimum width=2.2cm,minimum height=1.2cm,fill=blue!6] (net) at (3.1,0) {$\Net$};
  \node[below=1mm of net,font=\scriptsize,text width=2.6cm,align=center] {ReLU network\\(piecewise linear)};
  \draw[->] (x) -- (net);
  % matrix
  \begin{scope}[shift={(6.4,-1.05)}]
    \draw (0,0) rectangle (4.2,2.1);
    \draw (3.5,0) -- (3.5,2.1);
    \foreach \yy in {0.42,0.84,1.26,1.68} \draw[gray!50] (0,\yy) -- (4.2,\yy);
    \node[font=\scriptsize,align=center,text width=3.3cm] at (1.75,2.5) {$d$ columns: $\partial \Psi_c/\partial x_i \cdot x_i$\\[-1pt](gradient$\times$input, per class)};
    \node[font=\scriptsize,align=center,text width=1.6cm] at (3.85,2.55) {bias\\column};
    \node at (1.75,1.05) {$J(x)\,\mathrm{diag}(x)$};
    \node[font=\scriptsize] at (3.85,1.05) {$c(x)$};
    \draw[decorate,decoration={brace,mirror,amplitude=4pt}] (0,-0.12) -- (4.2,-0.12)
      node[midway,below=4pt,font=\scriptsize] {$\mathrm{M}(x)\in\mathbb{R}^{C\times(d+1)}$};
  \end{scope}
  \draw[->] (net) -- node[above,font=\scriptsize]{one matrix} (6.3,0);
  % row sums -> logits
  \begin{scope}[shift={(12.3,-1.05)}]
    \draw (0,0) rectangle (0.7,2.1);
    \foreach \yy in {0.42,0.84,1.26,1.68} \draw[gray!50] (0,\yy) -- (0.7,\yy);
    \node[font=\scriptsize,align=center,text width=2.4cm] at (0.35,2.5) {logits\\$\Netx$};
  \end{scope}
  \draw[->] (10.75,0) -- node[above,font=\scriptsize]{row sums}
            node[below,font=\scriptsize]{(exact)} (12.25,0);
\end{tikzpicture}
\caption{\textbf{The knowledge matrix encodes the germ of the network function.} At almost every
input, the network computes an affine map $x' \mapsto J x' + c$ on the surrounding linear
region --- its germ (Definition~\ref{def:germ}). The knowledge matrix arranges this germ as
one fixed-shape matrix
$\mathrm{M}(x)=[\,J(x)\,\mathrm{diag}(x) \mid c(x)\,]$: its first $d$ columns are per-class
gradient$\times$input contributions, its last column collects every bias term, and each row
sums \emph{exactly} to the corresponding logit. Because column $i$ is scaled by $x_i$, the
germ is \emph{recoverable} from the matrix at coordinates where $x_i\neq0$
(Remark~\ref{rem:zero-pixels}); at $x_i=0$ that column vanishes. Everything in this paper
follows from this picture: invariance (the germ does not change under reparameterisations that
preserve the function), completeness (the germ is recoverable from the matrix where the input
is nonzero), and the distance geometry
(the row-sum constraint splits any matrix change into a function-visible part and a
function-invisible remainder).}
\label{fig:germ}
\end{figure}
